% from ega0-19-fr.tex
\oldpage[0\textsubscript{IV-1}]{78}

\begin{remark}[19.1.15]
\label{0.19.1.15-fr}
La notion «~duale~» de celle d'homomorphisme formellement inversible à gauche
est celle d'homomorphisme \emph{formellement inversible à droite}~; un tel
$A$-homomorphisme continu $u:M\to N$ vérifie, par définition, la condition
suivante~: pour tout sous-module ouvert $W'$ de $N$, il existe un sous-module
ouvert $V'\subset u^{-1}(W')$ de $M$, un sous-module ouvert $W''\subset W'$ de
$N$ et un homomorphisme $h:N/W''\to M/V'$ tels que l'homomorphisme canonique
$N/W''\to N/W'$ se factorise en
\[
  N/W''\xrightarrow{h}M/V'\xrightarrow{u'}N/W',
\]
où $u'$ est déduit de $u$ par passage aux quotients. Cela entraîne que $u$ est
un épimorphisme formel~; s'il existe un $\widehat A$-homomorphisme continu
$r:\widehat N\to\widehat M$ tel que $\widehat u\circ r=1_{\widehat N}$, on
vérifie aussitôt que $u$ est formellement inversible à droite. Si les
conditions de (19.1.7) sont remplies, pour que $u$ soit formellement inversible
à droite, il faut et il suffit que les $u_\lambda$ soient inversibles à droite
(c'est-à-dire que le noyau de $u_\lambda$ est facteur direct de
$M/\mathfrak J_\lambda M$ et que $u_\lambda$ est un isomorphisme d'un
supplémentaire de $\operatorname{Ker}(u_\lambda)$ sur
$N/\mathfrak J_\lambda N$). Nous laissons au lecteur le soin d'énoncer et de
démontrer les propositions correspondant à (19.1.8) et (19.1.9) (dans
l'analogue de (19.1.8), il faut supposer $M$ séparé et complet et que $M_n$ est
un $(A/\mathfrak J_n)$-module projectif~; dans l'analogue de (19.1.9), il faut
supprimer l'hypothèse sur les $N/\mathfrak J_\lambda N$, mais supposer par
contre que, pour tout $\lambda$, $M/\mathfrak J_\lambda M$ est un
$(A/\mathfrak J_\lambda)$-module projectif).
\end{remark}

\subsection{Modules formellement projectifs}
\label{subsection:0.19.2-fr}

\begin{definition}[19.2.1]
\label{0.19.2.1-fr}
Soient $A$ un anneau topologique. On dit qu'un $A$-module topologique $M$ est
\emph{formellement projectif} s'il vérifie la condition suivante~: pour tout
idéal ouvert $\mathfrak J$ de $A$, tout couple de $(A/\mathfrak J)$-modules
(discrets) $P$, $Q$, tout $A$-homomorphisme surjectif $u:P\to Q$ et tout
$A$-homomorphisme continu $v:M\to Q$, il existe un $A$-homomorphisme continu
$w:M\to P$ tel que $v=u\circ w$.
\end{definition}

\begin{env}[19.2.2]
\label{0.19.2.2-fr}
Pour vérifier la condition de (19.2.1), on peut évidemment se borner (en
remplaçant $Q$ par $v(M)$ et $P$ par $u^{-1}(v(M))$) au cas où $v$ est
lui-même surjectif~; alors $Q$ est isomorphe à $M/V$, où $V$ est un
sous-module ouvert de $M$ tel que $\mathfrak JM\subset V$~; la condition de
(19.2.1) équivaut alors à dire que pour tout $(A/\mathfrak J)$-module discret
$P$ et tout $A$-homomorphisme surjectif $u:P\to M/V$, il existe un sous-module
ouvert $V'\subset V$ dans $M$ et un $A$-homomorphisme $w:M/V'\to P$, tels que
l'homomorphisme canonique $M/V'\to M/V$ se factorise en
\[
  M/V'\xrightarrow{w}P\xrightarrow{u}M/V.
\]
On notera qu'il suffit de vérifier cette condition lorsque $\mathfrak J$
parcourt un système fondamental de voisinages de 0 dans $A$ formé d'idéaux.
\end{env}

\begin{env}[19.2.3]
\label{0.19.2.3-fr}
Supposons qu'il existe un système fondamental $(\mathfrak J_\lambda)$ d'idéaux
ouverts dans $A$ et un système fondamental $(V_\lambda)$ de sous-modules
ouverts de $M$, ayant \emph{même} ensemble d'indices que
$(\mathfrak J_\lambda)$ et tel que, pour tout $\lambda$, $M/V_\lambda$ soit un
$(A/\mathfrak J_\lambda)$-module \emph{projectif}. Alors $M$ est formellement
projectif~: il suffit en effet, avec les notations de (19.2.2), de prendre
$\lambda$ tel que $\mathfrak J_\lambda\subset\mathfrak J$ et
$V_\lambda\subset V$~; comme $P$ est aussi un
$(A/\mathfrak J_\lambda)$-module, la factorisation de l'homomorphisme canonique
$M/V_\lambda\to M/V$ en $M/V_\lambda\to P\to M/V$ résulte de ce qu'il s'agit
d'un $(A/\mathfrak J_\lambda)$-homomorphisme et de ce que $M/V_\lambda$ est
supposé être un $(A/\mathfrak J_\lambda)$-module projectif.

Lorsque la condition plus stricte de ce numéro est satisfaite, on dit que $M$
est \emph{strictement formellement projectif}.
\end{env}

\begin{proposition}[19.2.4]
\label{0.19.2.4-fr}
Soient $A$ un anneau topologique, $M$ un $A$-module topologique dont la
topologie est déduite de celle de $A$. Pour que $M$ soit formellement
projectif, il faut et il suffit qu'il soit strictement formellement projectif.
\end{proposition}

